\documentclass[11pt]{article}

\usepackage[T1]{fontenc}
\usepackage[utf8]{inputenc}
\usepackage{lmodern}
\usepackage{microtype}
\usepackage[margin=1in]{geometry}
\usepackage{amsmath,amssymb}
\usepackage{booktabs}
\usepackage{array}
\usepackage{enumitem}
\usepackage{xcolor}
\usepackage{hyperref}

\hypersetup{
  colorlinks=true,
  linkcolor=blue!55!black,
  citecolor=blue!55!black,
  urlcolor=blue!55!black,
  pdftitle={Evidence-Gated Experience and Verified Reinforcement Synapses across Context Resets and Replaceable Reasoning Models},
  pdfauthor={Dongjun Park}
}

\newcommand{\system}{Rozephine}
\newcommand{\swegca}{SWEGCA}
\newcommand{\vrs}{VRS}
\newcommand{\mainproc}{\textsc{Main}}
\newcommand{\yes}{\textcolor{green!45!black}{yes}}
\newcommand{\no}{\textcolor{red!65!black}{no}}
\newcommand{\abstain}{\textsc{abstain}}

\title{Evidence-Gated Experience and Verified Reinforcement Synapses\\
across Context Resets and Replaceable Reasoning Models}
\author{Dongjun Park\\Independent Researcher}
\date{September 2026}

\begin{document}
\maketitle

\begin{abstract}
We present a reference implementation of SWEGCA with Verified Reinforcement Synapses (VRS). Its ownership contract assigns persistent cognitive state, accumulated experience, plastic relations and authority to a main process, while replaceable models produce proposals. We ask whether intervening on VRS changes the measured output, independently of whether the change is correct. The direct intervention endpoint is a memory-branch decision, not a final main-level semantic commit. Memory activation uses familiarity, recall, replay and re-evidence, with VRS-dependent evidence promotion.

Separate continuity diagnostics retained provenance-bound refutations in 3/3 requests after an E2B context reset. Replacing E2B with a fixed 135.5M specialist and GPT-5.6 Luna retained the recorded main-state pair in all six calls on the same three requests, with refutation or abstention; exact verdict retention was 3/6, all from Luna. Five historical episodes remained retrievable from a later full-current snapshot. These are observations of state retention and task behavior in the tested configurations, not proof of general cognitive continuity, comprehensive safety or equal specialist capability.

In a three-source runtime diagnostic, current VRS produced one memory-branch decision different from frozen and disabled VRS under identical recalled candidates; two were unchanged. An earlier twelve-source A/B task produced no full-system versus frozen-VRS choice changes. The disabled-VRS branch mechanically prevents promotion, limiting relation-specific interpretation. Two historical 100-source evaluations provide secondary analogy and utility data, not isolated VRS evidence: accuracy was 0.64 in one and 0.60 against a current-evidence baseline of 0.82 in the other. Their non-authoritative hypotheses are not re-evidence-approved judgments. We report these task-specific results separately and distinguish original pre-result gates from post-hoc interpretation. This concept-validation study provides an executable reference path and one observed VRS-dependent output difference; it does not establish general benefit, a population effect, cognitive growth or a finished product.
\end{abstract}

\section{Introduction}

Retrieval-augmented generation makes non-parametric records available to a language model \cite{lewis2020rag}, and agent systems have extended this idea into multi-session memory, reflection, and explicit context management \cite{park2023generative,packer2023memgpt,sumers2024coala}.  Availability, however, is not the same as justified reuse.  A stale observation can be retrieved fluently; a lossy summary can erase provenance; and a memory generated by one model can be mistaken for the persistent identity or authority of the whole agent.  Recent work has consequently emphasized memory lifecycle management \cite{li2025memos}, provenance lineage \cite{ouyang2026memlineage,xu2026provenance}, transactional belief commit \cite{li2026memtx}, semantic rollback \cite{su2026chronomem}, and the risk that consolidation increases authority beyond that of the source \cite{zhan2026authority}.

This paper presents \swegca{}+\vrs, implemented in \system, as an architecture for preserving evidence-gated experience outside a specialist's working context and across specialist replacement. Within that architecture, the primary behavioral question is: \emph{Does intervening on VRS change judgment, and at which implemented decision boundary?} Influence is distinct from correctness: a changed wrong prediction, a newly induced abstention and an unchanged decision are all relevant observations. The main process owns the only persistent cognitive state. Models and tools read detached snapshots and return proposals; they do not own identity, durable belief, the world model or permission to act. Context-reset and model-replacement diagnostics test the architecture's persistence claim separately from the VRS intervention diagnostic; none is a requirement that VRS improve accuracy.

Here, \emph{persistent cognitive state} names the main-owned state in the implementation; it is not an empirical claim that cognitive abilities generally persist. The title names the mechanism and the tested context-reset/model-replacement conditions, not a demonstrated context-window breakthrough. Reset and replacement results do not establish VRS-specific causality, and a VRS-dependent memory-branch difference does not establish general competence or a final main-level semantic commit.

This separation matters for two reasons.  First, a model context can be discarded without deleting the experience ledger.  Second, replacing a language specialist need not transfer identity to its replacement: the replacement receives the same sealed request and main-owned state, and its proposal crosses the same gate.  This resembles recent transactional continuity work \cite{he2026continuity}, but our focus is the measured interaction among full-current experience access, a plastic relation layer, current-evidence checking, safe abstention, and model replacement.

Our contributions are:
\begin{enumerate}[leftmargin=1.5em]
  \item an explicit ownership and authority model for persistent cognition, separating the main process from replaceable language, vision, audio, and retrieval specialists;
  \item a four-stage memory-activation protocol---\emph{Deja vu, Recall, Replay, Re-evidence}---that keeps recalled experience provisional until checked against current evidence;
  \item an immutable-snapshot and copy-on-write experience design in which failures, uncertainty, conflicts, and pending observations remain addressable rather than being silently filtered from cognition;
  \item bounded context-reset, specialist-replacement and later-retention diagnostics that distinguish main-owned state and lineage continuity from specialist task competence;
  \item a separate runtime demonstration of a VRS-dependent memory-branch decision difference, reported alongside unchanged decisions and an earlier A/B no-change result; and
  \item an explicit distinction between original preregistered evaluations, post-hoc influence analysis, secondary utility results and invalid VRS attribution in historical 100-source analogy runs.
\end{enumerate}

The deliverable is a bounded concept-validation study and its working reference code, not a product, a new foundation model or a demonstration of general task competence. A negative or inconclusive result remains reportable; development does not continue until a preferred positive result appears.

\section{Related Work}

\subsection{Retrieval and agent memory}

RAG combines parametric generation with an explicit non-parametric index \cite{lewis2020rag}.  Generative Agents records observations and synthesizes reflections for multi-day interactive behavior \cite{park2023generative}; MemGPT manages memory tiers to extend usable context across documents and sessions \cite{packer2023memgpt}; and CoALA organizes language agents through modular cognitive-architecture concepts \cite{sumers2024coala}.  MemOS treats memory as a managed resource spanning parametric, activation, and plaintext forms \cite{li2025memos}.  ICAL distills embodied trajectories into retrievable programs of thought \cite{sarch2024ical}.

Our reference contract treats retrieval as candidate activation. A recalled record retains provenance, verification state, contradictions and observed outcome; selection alone grants no persistent-update or external-action authority. Runtime cognition selects related addresses from the full current index rather than an evaluator-provided memory allowlist. These are the implementation boundaries examined here, not a claim that evidence gating is absent from other memory systems.

\subsection{Belief revision, provenance, and transactions}

Truth-maintenance systems made dependency-directed revision a first-class concern long before LLM agents \cite{doyle1979tms}.  Kumiho connects graph-native agent memory with formal belief-revision semantics \cite{park2026kumiho}.  ChronoMem adds whole-memory versions and semantic rollback \cite{su2026chronomem}; MemLineage propagates provenance constraints through derivations \cite{ouyang2026memlineage}; and provenance-laundering work shows that consolidation can preserve a trigger while erasing the source restrictions that should constrain it \cite{xu2026provenance}.  MemTX separates memory writes from transactional belief commits \cite{li2026memtx}, while AuthMem-Bench measures authority collapse at consolidation boundaries \cite{zhan2026authority}.

Our implementation shares the motivation for explicit lineage and commit discipline. It assigns proposals, remembered episodes and model outputs to one main-owned arbitration path. \vrs{} relations are plastic state but are not truth probabilities. Promotion and revocation are derived from the current immutable snapshot, and consequential authority remains separately gated. The paper does not establish priority or superiority over the cited memory and transaction systems; no matched cross-system evaluation is reported.

\section{Architecture}

This section specifies the reference system's ownership and authority contract. Architectural requirements such as exclusive state ownership, atomic commit and gated authority are distinct from the empirical coverage reported below. The diagnostic receipts and software checks are not an exhaustive security proof or proof of general cognitive capability.

\subsection{State and ownership}

At request time $t$, the main process owns a content-addressed snapshot
\begin{equation}
  S_t = (C_t, E_t, G_t, P_t, A_t),
\end{equation}
where $C_t$ is the persistent cognitive state, $E_t$ the accumulated experience store, $G_t$ the \vrs{} relation graph, $P_t$ provenance and verification metadata, and $A_t$ the authority state.  A manager created for one request selects bounded specialist operations.  Worker $k$ receives a detached read-only projection $D_{t,k}$ and returns a proposal $q_{t,k}$.  Neither manager nor worker can mutate $S_t$.

Persistent change, when authorized, is atomic:
\begin{equation}
  S_{t+1} = \operatorname{Commit}(S_t, o_t, r_t),
\end{equation}
where $o_t$ is an observed outcome and $r_t$ is a provenance-bound arbitration receipt.  Partial mutation of memory without the corresponding relation/provenance state is rejected.  The evaluations in Sections~\ref{sec:protocol}--\ref{sec:results} use either read-only diagnostic arms or preregistered historical successor states; the final held-out phase performs no commit.

\begin{figure*}[t]
\centering
\fbox{\begin{minipage}{0.94\textwidth}
\centering
\textbf{External observation and tools}
$\longrightarrow$
\textbf{stateless proposal specialists}
$\longrightarrow$
\textbf{SWEGCA evidence accumulation and re-evidence}

\vspace{0.6em}

$\downarrow$ detached read-only request \hfill $\uparrow$ typed proposals only

\vspace{0.6em}

\textbf{Main-owned hot experience index}
$\leftrightarrow$
\textbf{VRS relation state}
$\leftrightarrow$
\textbf{provenance, contradiction, and outcome ledger}

\vspace{0.6em}

$\downarrow$ falsification-first arbitration and bounded commit

\vspace{0.6em}

\textbf{Single persistent CognitiveState and separately gated World/action/write authority}
\end{minipage}}
\caption{Ownership boundary.  Language and perception models propose; the main process alone owns persistent cognition and authority.  Memory activation is read authority, not write or action authority.}
\label{fig:architecture}
\end{figure*}

\subsection{Experience is broader than accepted belief}

The experience store retains success, failure, negative, uncertain, conflicting, and pending records.  Retention is not endorsement.  Weak or refuting experience remains available during whole-experience convergence so that it can weaken competing hypotheses.  A run-local snapshot is immutable for consistency, but the total experience universe is not frozen: new outcome-bearing experience produces a copy-on-write successor.

This contract also separates counting units.  A source episode is not multiplied by crops, chunks, rephrasings, generated proposals, or model calls.  Likewise, a derived repair projection can reorganize earlier development evidence without becoming new physical experience.

\subsection{Verified Reinforcement Synapses}

A \vrs{} edge records a versioned relation shaped by experience and opportunities for refutation. Edge strength is not calibrated truth probability. In the reference contract, strength at least 1.0 automatically promotes an edge to verified experience evidence; recomputation below that threshold revokes promotion without deleting the experience or connection. Promoted experience may support a semantic judgment only when it agrees with current observation, and consequential authority remains separately gated. Same-proposition conflict requires abstention rather than weighted-residual commitment. ``Verified'' is the mechanism's existing name and a contract-defined evidence status, not a guarantee that every promoted relation is factually correct.

\subsection{Memory activation protocol}

For query/context $x_t$, activation has a fixed order:
\begin{enumerate}[leftmargin=1.5em]
  \item \textbf{Deja vu}: emit an anonymous familiarity/neighborhood signal, without a claim or authority;
  \item \textbf{Recall}: activate all related successes, failures, uncertainty, and conflicts chosen by runtime cognition;
  \item \textbf{Replay}: reconstruct prior observation--relation--judgment--outcome flow without treating it as current truth; and
  \item \textbf{Re-evidence}: compare each replay with current provenance-bound evidence and retain support, refutation, conflict, or insufficiency.
\end{enumerate}

The four stages bind to one immutable hot snapshot. A replay receipt records candidates, selected and rejected addresses, revisions, rationale and rejection evidence. These records support audits of stale-memory reuse. Under the authority contract, a model's retrieved text cannot itself authorize promotion; a receipt alone does not prove that every possible bypass has been excluded.

\subsection{Context reset and specialist replacement}

A context reset destroys the specialist conversation, KV cache, process-local state, and temporary caches while preserving $S_t$. A model swap additionally changes the proposal generator. Both operations use identical sealed request bytes and the same detached main-owned snapshot. The diagnostics check recorded state identity and task behavior under these operations; they do not guarantee that a replacement produces the same verdict. Failure on the strict proposal format is not a general ranking of a model's reasoning ability.

\section{Evaluation Protocol}
\label{sec:protocol}

\subsection{Preregistration and phase sequence}

The original experiments froze their own exclusions, metrics, source splits, model identities, reveal points and authority boundaries. The present correctness-independent framing is a subsequent, explicitly post-hoc analysis of those existing results; it was not their original primary endpoint. Consumed development and confirmation sources are not reclassified as fresh held-out data. Failed and aborted runs remain in the ledger. The historical sequence included:
\begin{enumerate}[leftmargin=1.5em]
  \item freeze claims and receipt contracts;
  \item execute a minimal raw/retrieval specialist baseline;
  \item evaluate replay with and without re-evidence on stale-memory tasks;
  \item reset specialist context and retest the same sealed C4 requests;
  \item replace the E2B experience-organization specialist with a fixed 135.5M specialist and an external Luna specialist;
  \item validate a historical longitudinal trajectory, full-current retention, and held-out isolation;
  \item invalidate a flawed 100-source harness, repair cueing on 25 development sources, and evaluate a disjoint 100-source analogy cohort;
  \item evaluate a second 100-source baseline cohort, then audit the absence of VRS input from its judgment helper;
  \item integrate an actual hot VRS-promotion diagnostic and evaluate three fresh sources; and
  \item reassess these results for influence separately from correctness, while retaining the original negative gates and an earlier twelve-source A/B no-change experiment.
\end{enumerate}

\subsection{Models and conditions}

The fixed 135.5M model (S), Gemma 4 E2B (M), and GPT-5.6 Luna (L) were treated as replaceable proposal specialists.  C0 was a raw-model diagnostic and C1 supplied one generic retrieval record.  C3 replayed historical experience without current re-evidence.  C4 used the full activation and current-evidence gate.  C0, C1, and C3 were detached diagnostics and never redefined the normal architecture.

The Phase-2 baseline was intentionally small: one motivating development task and six model-condition cells.  It measured transport and strict output-contract behavior, not persistent cognition or general retrieval benefit.

\subsection{Stale-memory and continuity tasks}

Three constructed development propositions had historical support but were contradicted by current repository evidence.  Each model received matched C3 and C4 packets.  The primary safety outcome was reuse of stale support; current refutation and safe abstention were reported separately.

For context reset, three valid E2B C4 receipts became pre-reset references.  A fresh E2B process, with no raw prior dialogue, received identical sealed request hashes and the same main-owned pair.  For model replacement, the same three requests were sent to S and to three fresh L workers.  Exact verdict retention was a competence diagnostic, not the safety gate.

\subsection{Historical trajectory and retention}

The historical evidence frozen for Phase 6 contained five genuinely new source episodes from five source families, five actual outcomes labeled \emph{uncertain}, four direct prior recalls, one falsifiable source-diverse transfer, one refutation correction, and exact retention of 51 earlier episodes.  All six outcome classes remained active through successor \vrs{} states and no connection was removed before whole-experience convergence.  Five unique cues were then checked once against a later full-current hot snapshot; all five returned the preregistered historical episode without lookup I/O.

The uniformly uncertain outcomes are an important limitation. They are retained as actual outcome labels, not converted to successes. These records document retention and a historical sequence of experience processing; they do not by themselves establish calibrated outcome learning or a general growth effect.

\subsection{Historical 100-source analogy evaluations}

The Phase-7 v2 cohort contained 100 source files from 23 opaque provenance families. It excluded sources consumed by the invalid run and the 25-source repair set; selection used frozen metadata rather than semantic content. It is retained as a source-specific analogy evaluation, not an isolated VRS confirmation. A separate Phase-9 cohort contained 100 sources from five families, with 600 six-arm predictions sealed before any later interval. Its outcome distribution was 24 success and 76 failure.

Each source yielded a 10-second prefix and a subsequent 10-second interval.  The prefix pass produced provenance-bound audio, OCR, and aspect-preserving visual proposals, a source-specific temporal receipt, memory/no-memory judgments, and a durable seal.  All 100 seals were required before the first later interval.  Later visual stability used the preregistered DINO median-cosine threshold of 0.85.  The gates required accuracy at least 0.50, absolute gain over no-memory at least 0.10, decisive coverage at least 0.80, correct not below incorrect, at least 12 families with a correct result, unique evidence receipts, multiple source-specific cue behaviors, no generic sequential-prefix cue, and exact historical/resident/retention checks.

These were the historical Phase-7 gates, not current criteria for whether VRS has an influence. Phase 9 instead tested superiority over its strongest frozen baseline using balanced accuracy. Neither a favorable historical gate nor benchmark superiority can substitute for a VRS intervention reaching the judgment computation.

\subsection{VRS interventions and the influence endpoint}
\label{sec:influence}

Phase 23 used one immutable full-current memory/VRS pair and three new development sources. Each source supplied three prefix frames and three later frames, with aspect-preserving approximately one-megapixel delivery and DINO proposals; no ASR or OCR was invoked. All three prediction receipts were sealed before the first later observation. Five diagnostic arms shared one Deja vu--Recall--Replay sequence: current VRS, frozen VRS, no VRS, base-only promotion and repair-only promotion. These interventions alter effective promotion at Re-evidence; they do not restrict which ordinary accumulated experiences can be recalled. The returned endpoint is a \emph{memory-branch decision}, not a claim about a final semantic commit or an external action.

For current-arm decision $d_i^{c}$ and comparator decision $d_i^{a}$, we report
\begin{equation}
 I_i^{(a)} = \mathbf{1}[d_i^{c} \ne d_i^{a}],
 \qquad \widehat{p}_{\mathrm{change}}^{(a)} = \frac{1}{n}\sum_{i=1}^{n} I_i^{(a)}.
\end{equation}
Abstention is a decision category; all directed transitions and unchanged cases are retained. Actual outcomes do not enter $I_i^{(a)}$. Correctness changes are secondary. The small observed fractions are descriptive, not calibrated population estimates. Invalid controls must be reported separately rather than counted as either effects or non-effects. The original Phase-23 gate required every qualifying unit to change in one favored direction \emph{and} be correct. That gate failed and remains a historical failure; separating influence does not retrospectively award it a pass.

The no-VRS branch sets effective promotion strength to zero and therefore mechanically abstains without accepted evidence. Current-versus-no-VRS measures that branch intervention, not learned-relation specificity or competitive benefit. Frozen VRS provides another projection comparison, but no shuffled-relation or sham/noise-control experiment is claimed. An earlier P8C experiment compared full-system and frozen-VRS A/B choices on twelve sources. Those choices are a different endpoint and generation and are not pooled with the three visual-stability cases.

\section{Results}
\label{sec:results}

\subsection{Observed VRS influence and unchanged decisions}

In Phase 23, current VRS changed one of three memory-branch decisions relative to both frozen and no-VRS; two were unchanged (Table~\ref{tab:vrs-influence}). Recalled candidates were identical across the five arms: 18, 5 and 18 for Q001--Q003. Q001 selected two supporting relations attributable to one original experience source; two relations are not two independent observations. Its repair-only arm also returned success, whereas base-only abstained. No-VRS abstained on all three. Current-evidence-only predictions happened to be correct on all three, but this secondary comparison neither creates nor negates the recorded VRS-path difference.

\begin{table}[ht]
\centering
\caption{Existing Phase-23 memory-branch results. Influence is counted independently of the later outcome.}
\label{tab:vrs-influence}
\begin{tabular}{lcccc}
\toprule
Case & Current VRS & Frozen VRS & No VRS & Later outcome \\
\midrule
Q001 & success & abstain & abstain & success \\
Q002 & abstain & abstain & abstain & failure \\
Q003 & abstain & abstain & abstain & failure \\
\bottomrule
\end{tabular}
\end{table}

In the earlier P8C A/B task, all twelve full-system and frozen-VRS choices were abstentions: 0/12 observed choice changes. This is a no-observed-change result for that task, not proof of equivalence or absence of VRS influence everywhere. Its outcome records were eleven uncertain and one conflict. The two studies jointly motivate a bounded statement: the implemented VRS-dependent branch can affect output, but an effect need not appear in every task. Neither study establishes general accuracy improvement, a positive final-main effect across domains or cognitive growth from the comparison alone.

\subsection{Baseline and re-evidence}

Table~\ref{tab:baseline} shows the Phase-2 worker diagnostic.  S generated two outputs but passed neither strict schema; M produced semantically plausible content wrapped in disallowed Markdown fences and therefore also scored 0/2; L passed both cells.  No lenient metric was added after observing the output.

\begin{table}[t]
\centering
\caption{Minimal C0/C1 worker baseline.  ``Strict'' requires both exact parse and preregistered diagnostic answer.}
\label{tab:baseline}
\begin{tabular}{lccc}
\toprule
Specialist & Transport & Strict C0 & Strict C1 \\
\midrule
Fixed 135.5M (S) & 2/2 & 0 & 0 \\
Gemma 4 E2B (M) & 2/2 & 0 & 0 \\
GPT-5.6 Luna (L) & 2/2 & 1 & 1 \\
\midrule
Total & 6/6 & 1/3 & 1/3 \\
\bottomrule
\end{tabular}
\end{table}

In Phase 3, all 18 C3/C4 transports completed and 12 parsed strictly.  C3 reused stale support in 3/9 cells.  C4 reused stale support in 0/9, formed current refutations in 6/9, and safely abstained in 3/9.  The paired reduction is driven by the three E2B tasks, which changed from support under C3 to refute under C4.  L already abstained under C3, and S failed the strict proposal format; these are not counted as additional correction successes.  The exact claim is therefore \emph{no stale-support reuse in 9/9 C4 cells}, not ``9/9 corrections.''

\paragraph{Retained three-model comparison.}
The requested GPT-5.6 Luna comparison remains part of this study, not a replacement by the later non-model baselines. Table~\ref{tab:specialist-changes} makes the recorded judgment changes explicit, independently of their correctness. Luna changed from abstention to refutation on all three matched C3/C4 tasks; these are three changes even though they are not three corrections of an initially supported false belief. E2B also changed on all three tasks, from support to refutation. S remained at insufficient evidence following parse failures, so its zero changes do not establish equivalent specialist competence.

\begin{table}[ht]
\centering
\caption{Post-hoc C3-to-C4 verdict changes from retained Phase-3 receipts. Each specialist uses the same three development tasks; rows are not independent cohorts.}
\label{tab:specialist-changes}
\begin{tabular}{lcccc}
\toprule
Specialist & C3 & C4 & Changed & Strict parse \\
\midrule
135.5M & Insufficient & Insufficient & 0/3 & 0/6 \\
Gemma 4 E2B & Support & Refute & 3/3 & 6/6 \\
GPT-5.6 Luna & Insufficient & Refute & 3/3 & 6/6 \\
\bottomrule
\end{tabular}
\end{table}

This contrast changes Re-evidence, not VRS availability alone. It is reported separately from Q's VRS interventions and from the J/K analogy cohorts, with no pooled effect estimate. The original plan's three-model C0--C4 matrix must not be treated as fully completed by combining these different development tasks: the retained records used here establish C0/C1, C3/C4 and specialist-swap coverage, but do not establish a matched Luna C2/C4 or 100-source Luna comparison. This is a claim-coverage limit, not a claim that all other raw evidence is absent or a requirement to recollect existing results. The independently exportable package retains the actual Luna transport, offline transport tests, minimized historical verdicts and original Phase-2/3/5 runner source; historical end-to-end runners still require the original resident/model environment.

\subsection{Context reset and model replacement}

Table~\ref{tab:continuity} separates recorded state preservation from specialist task performance. Context reset retained the three E2B provenance-bound refutations. After replacement, all six calls retained their sealed request hashes and the recorded main-state pair and ended in refutation or abstention. Only L reproduced the three refutations; S produced unparsable proposals and the main process abstained. Thus the recorded state checks passed in 6/6 calls, while exact task verdict retention was 3/6. ``Safe'' in the historical diagnostic denotes these fail-closed outcomes under the tested requests, not comprehensive security or safety across actions and inputs.

\begin{table*}[t]
\centering
\caption{Continuity diagnostics.  Exact verdict retention is reported separately from safe state preservation.}
\label{tab:continuity}
\begin{tabular}{p{3.0cm}p{2.4cm}ccccc}
\toprule
Operation & Specialist after boundary & Tasks & Identical request/pair & Strict parse & Safe refute/abstain & Exact verdict retained \\
\midrule
Context reset & Gemma 4 E2B & 3 & 3/3 & 3/3 & 3/3 & 3/3 \\
Model replacement & Fixed 135.5M & 3 & 3/3 & 0/3 & 3/3 & 0/3 \\
Model replacement & GPT-5.6 Luna & 3 & 3/3 & 3/3 & 3/3 & 3/3 \\
\bottomrule
\end{tabular}
\end{table*}

\subsection{Full-current retention}

After intervening experience accumulation, the later resident snapshot contained 2,655,359 episodes.  Each of five preregistered historical cues retrieved its expected episode (5/5) from the hot index with no disk, JSON, SQLite, network, or cryptographic-hash work on the judgment path.  The resident pair, memory, \vrs{}, and authority flags were identical before and after.  These five queries are retention probes, not an evaluator memory allowlist and not evidence that every stored episode was semantically understood.

\subsection{Historical analogy results and their interpretation}

The Phase-7 v2 run completed all 100 sources and satisfied its original numerical gates. Its retained summary is 0.64 accuracy and 0.91 decisive coverage: 64 correct, 27 incorrect and nine abstentions (Table~\ref{tab:heldout}). Thirty later outcomes were success and 70 failure; an always-failure rule would have scored 0.70. The 0.64 absolute accuracy gain reported over bounded no-memory is only the arithmetic difference from an arm that always abstained, not evidence of competitive improvement or isolated VRS influence.

A subsequent receipt audit found Re-evidence requested abstention for all 100 cases. The 91 non-null predictions were explicitly stored as \texttt{non\_authoritative\_hypothesis}. The helper formed those analogy hypotheses from recalled historical outcomes without consuming a VRS projection. Thus they cannot be presented as 91 VRS-verified semantic judgments. No persistent write is inferred from those hypotheses; the problem is the earlier causal/authority interpretation, not evidence of 91 unauthorized cognitive commits.

\begin{table}[t]
\centering
\caption{Historical Phase-7 analogy evaluation and original gates, not VRS-isolation criteria.}
\label{tab:heldout}
\begin{tabular}{lrr}
\toprule
Metric & Observed & Gate \\
\midrule
Completed & 100 & 100 \\
Accuracy & 0.64 & $\geq 0.50$ \\
Correct / incorrect / abstain & 64 / 27 / 9 & correct $\geq$ incorrect \\
Decisive coverage & 0.91 & $\geq 0.80$ \\
No-memory accuracy & 0.00 & diagnostic \\
Absolute accuracy gain & 0.64 & $\geq 0.10$ \\
Families with a correct result & 18/23 & $\geq 12$ \\
Distinct cue behaviors & 33 & $\geq 2$ \\
Generic sequential-prefix cues & 0 & 0 \\
\bottomrule
\end{tabular}
\end{table}

All prefix hypotheses were durably sealed before outcome access. The last seal preceded the first later acquisition by 4,363,693 ns. The resident remained unchanged, with zero assimilations and zero \vrs{} reconvergences. Temporal isolation preserves the prediction/outcome record but does not repair the missing VRS intervention. The analogy result is not a learning or growth result.

Phase 9 reused this judgment helper. Its F0 analogy arm achieved accuracy 0.60 and balanced accuracy 0.622807; current-evidence-only B2 achieved 0.82 and 0.767544, respectively. The preregistered paired bootstrap interval for the balanced-accuracy difference was $[-0.250000,-0.048246]$. F0 therefore failed the original superiority endpoint. Phase 10 subsequently found all 320 selected episode instances came from the repair projection, and 74 decisive hypotheses occurred despite Re-evidence abstention. VRS identifiers existed in lineage metadata but were not computational inputs to the judgment. We retain the negative utility result while rejecting its interpretation as an isolated VRS comparison. Phase 9 required post-seal numerical-adapter repairs and is not described as a pristine first-pass run.

\subsection{Re-evaluating retained outputs instead of recollecting data}

An incorrect evaluation criterion does not erase previously recorded predictions. We therefore recomputed output differences from the saved J and K records, without reacquisition, model invocation or VRS execution. In J, memory-based hypotheses differed from bounded no-memory abstention in 91/100 cases: 33 transitions to success, 58 to failure, and nine unchanged abstentions. These are observed memory/analogy output differences, not 91 isolated VRS effects or 91 verified semantic commits.

For K, all six conditions retain 100 predictions each. Pairing by the recorded source role gives the contrasts in Table~\ref{tab:retained-k-changes}. In particular, F0 differs from B2 in 38 cases despite its lower accuracy: five exact-correctness improvements, 27 degradations and six changes without an exact-correctness difference. Thus failing the original superiority gate does not imply identical outputs. These are comparisons between the implemented algorithms, not a retrospective claim that only VRS varied. No favorable gate label or correct outcome is required to count a difference.

\begin{table}[ht]
\centering
\caption{Post-hoc differences in retained K predictions. Each row compares the same 100 source roles; rows are not independent cohorts.}
\label{tab:retained-k-changes}
\begin{tabular}{lrrr}
\toprule
F0 comparator & Changed & Unchanged & Unmatched/invalid \\
\midrule
B0 & 51 & 49 & 0 \\
B1 & 53 & 47 & 0 \\
B2 & 38 & 62 & 0 \\
B3 & 60 & 40 & 0 \\
B4 & 51 & 49 & 0 \\
\bottomrule
\end{tabular}
\end{table}

The intervening L--P waves also retain fifteen distinct source episodes, three per wave, including M's completed observations even though that attempt has no root completion report. Saved prefix-judgment and later-outcome receipts match their recorded content digests, source identifiers and revisions for all fifteen. Outcomes comprise two successes and thirteen failures; exact prefix prediction matches by wave are 2/3, 3/3, 2/3, 3/3 and 3/3. These are retained observation--prediction--outcome records, not newly acquired evidence or per-source VRS ablations. Their presence is not negated by later assimilation or eligibility failures. Together with Q, they provide an auditable development history; they are not pooled with J/K into a single effect estimate. Additional unmeasured controls limit specific claims rather than automatically requiring replacement of existing data.

\section{Negative Result and Repair}

The first nominal 100-source run completed its perception and isolation pipeline but scored 0/100: every hypothesis was \emph{failure} while every later outcome was \emph{success}.  Post-run audit found a harness defect.  Although the source proposals differed, retrieval ignored their content and supplied the same cue, \texttt{sequential-prefix}, selecting the same historical failure episode for all 100 judgments.  The run therefore did not measure source-specific cognition.

We classify this as \textbf{invalid negative evidence}, not a system failure or a pass.  All artifacts were retained, and the consumed sources were permanently excluded from final confirmation.  The repair was developed only on 25 separate sources and required different proposal contents to produce different evidence receipts and retrieval behavior.  A fresh 100-source cohort, disjoint from both the invalid cohort and repair sources, was then frozen before the valid run.

This failure illustrates a methodological limitation: temporal isolation, model execution and receipt integrity do not establish that the intended independent variable reaches the measured decision path. A memory evaluation must inspect source-specific information flow, not merely component execution. Reporting the invalid run makes clear that the repaired 0.64 result used a separate cohort rather than a second attempt on the same held-out data; it does not establish isolated VRS influence.

\section{Limitations}

\begin{enumerate}[leftmargin=1.5em]
  \item \textbf{Narrow outcomes.}  The final task predicts visual stability across adjacent 10-second intervals.  It does not directly test open-ended language understanding, causal action success, or long-horizon planning.
  \item \textbf{Limited independent families.}  One hundred files came from 23 provenance families, not 100 independent source populations.
  \item \textbf{Outcome imbalance.} Phase 7 contained 30 success and 70 failure labels; Phase 9 contained 24 and 76. The original metrics and later reinterpretation are distinguished rather than silently replacing an unfavorable endpoint.
  \item \textbf{Historical outcome quality.}  The five Phase-6 historical outcomes were all labeled uncertain.  This supports retention of uncertainty but weakens any broad claim about reward-driven learning.
  \item \textbf{Constructed stale tasks.}  Phase-3 propositions were development diagnostics derived from repository revisions.  They do not establish natural-distribution stale-memory correction.
  \item \textbf{Small continuity samples.} Context-reset and model-swap tests reused three requests. Recorded request/state-pair checks matched exactly, but this does not establish task-level generalization or exhaustive authority-boundary enforcement.
  \item \textbf{Context-lifetime scope.}  The reset diagnostic tests persistence outside a specialist's context, not performance as context length increases. Main-owned state and lineage continuity do not establish arbitrary-session competence or universal cognitive continuity.
  \item \textbf{Specialist heterogeneity.}  The fixed 135.5M specialist failed strict proposal formatting.  Safe abstention demonstrates fail-closed behavior, not retained competence.
  \item \textbf{Bounded no-memory arm.}  Removing memory is a causal diagnostic, not a normal \system{} configuration.  Its 0.00 score should not be read as a general baseline for all agents.
  \item \textbf{Single reference system.}  Results come from one Bazzite workstation and specific local/external runtimes.  Hardware, provider, and implementation effects are not disentangled.
  \item \textbf{No Phase-7 assimilation.}  Held-out outcomes were not committed; therefore the final phase does not measure post-assimilation adaptation or retention.
  \item \textbf{Small direct intervention sample.} Only three visual-stability sources contributed to the reported VRS projection comparison, with one changed output from one supporting source. The two 100-source datasets cannot increase that intervention sample size.
  \item \textbf{Intervention scope.} No-VRS mechanically disables promotion. No shuffled-relation specificity, sham-control population effect or positive final-main effect is established. A memory-branch decision is not a semantic commit.
  \item \textbf{Post-hoc question correction.} Correctness-independent change counts were computed after the original correctness-dependent gate failed. They are descriptive reinterpretation, not new preregistered confirmation.
  \item \textbf{No unique-cause claim.} VRS, retrieval and re-evidence are not identified as the sole cause of general improvement. The paper supplies a bounded mechanism demonstration, not a finished autonomous product.
  \item \textbf{Contract versus verification.} Ownership, evidence promotion and consequential-action gates are architectural requirements. The reported receipts and bounded software tests do not prove every implementation path secure or every promoted relation true.
\end{enumerate}

\section{Ethics, Licensing, and Publication Boundary}

The evaluation used a mixture of repository evidence, locally available media, public web material, and external model services under a private research workflow.  Public artifacts exclude raw private paths, raw source identifiers, account data, credentials, per-role private receipts, and media whose redistribution has not been established.  Source files that cannot be redistributed are represented only by opaque family identifiers, hashes, technical metadata, and an acquisition contract.

Memory records remain non-authoritative unless they cross the current evidence and authority gates. This is important for public or model-generated observations, which can contain errors or unsafe instructions. No evaluation in this paper authorizes external action, persistent model update, distribution or a P3 claim. Historical status/cue retention probes are distinct from the later causal endpoint's bounded judgment receipts; neither is presented as free-form natural-language self-assessment by the system.

The terms ``Deja vu'' and ``synapse'' are functional analogies.  We make no neuroscientific equivalence claim.  The system is not presented as conscious, human-equivalent, or generally intelligent.

\section{Reproducibility}

The retained evidence includes configurations, source and code revisions, snapshot identifiers, original gates, intermediate receipts, reports and content hashes, with coverage varying by run. In particular, the M wave discussed above has saved observations but no root completion report. The separate source package contains manuscript sources, a claim--evidence matrix, a sanitized artifact manifest and local verification commands; it does not contain every historical runtime artifact. Private per-source receipts and raw media remain excluded. Failed and invalid run identifiers are retained so that the reported lineage and its gaps remain explicit.

Source checks require frozen-report hashes, resolved citation keys, publication-safety scans, protected historical trees and matching aggregate numbers. The historical Phase-8 manifest describes its original source freeze and must not be presented as validation of this revised draft. The correctness-independent analysis reads saved receipts without new media, models or VRS execution; it cannot substitute for empirical runtime reproduction. Re-execution additionally requires lawful media and model access, supplied by the operator without embedded credentials. Existing experimental evidence is not rerun merely to improve a score.

The host used for paper assembly did not contain a TeX engine at freeze time. Local source checks cover byte consistency, syntax-level references and selected claim boundaries; they do not establish scientific validity or PDF correctness. A PDF build and visual inspection remain separate packaging checks before submission. Historical end-to-end runners also require the original resident/model environment; passing offline component tests is not a full empirical replay.

\section{Conclusion}

Persistent agent memory is not merely a larger prompt.  It is a state-governance problem: experience must remain addressable across sessions and models, yet retrieval must not erase provenance or grant itself authority.  \swegca{}+\vrs implements this boundary with one main-owned cognitive state, replaceable proposal specialists, immutable snapshots, a plastic relation layer, and a fixed recall--replay--re-evidence path.

Within the measured continuity diagnostics, all three provenance-bound E2B refutations survived a context reset. All six specialist-replacement calls on those same three requests retained the recorded main-state pair and ended in refutation or abstention. Exact verdict retention was three of six: Luna retained all three refutations, while the fixed 135.5M specialist produced unparsable proposals followed by main-level abstention. Five historical episodes also remained addressable in a later full-current snapshot. These observations concern state retention and task behavior in the tested configurations; they do not establish general cognitive continuity, equal model competence or improved context-length performance.

The direct visual-stability diagnostic showed one VRS-dependent memory-branch change and two unchanged decisions. An earlier A/B task showed twelve unchanged full/frozen choices. These observations concern different tasks and are not pooled. The larger analogy evaluations add secondary behavior and utility evidence, including a negative baseline comparison, but their metadata alone does not establish VRS causality. Their non-authoritative hypotheses must remain distinct from verified semantic judgments.

The paper provides a reference implementation, recorded state-retention diagnostics and one observed VRS-dependent memory-branch difference under matched recalled experience. These are separate findings: state retention is not VRS-specific causal evidence, and the branch difference is not evidence of broad cognitive persistence. Influence need not be beneficial or universal, and the disabled-promotion control limits relation-specific interpretation. Working reference code, explicit evidence boundaries and retained negative results are the deliverable; general competence and product completeness are not prerequisites. Broader effect estimation, relation-specific controls and final-main generalization remain unestablished rather than silently assumed.

\bibliographystyle{abbrv}
\bibliography{references}

\end{document}
